\mainlanguage[de]

\setuplayout
  [backspace=2.0cm,
   height=middle,
   width=middle,
   leftmargin=1.8cm,
   footerdistance=0.0cm,
   rightmargindistance=0.2cm,
   rightmargin=1.8cm,
   bottom=0.5cm,
   leftmargindistance=0.2cm,
   header=2.1cm,
   headerdistance=0.5cm,
   footer=0.5cm,
   top=0cm,
   topspace=0.5cm,
   bottomspace=0.0cm]

\setupformulas[align=right]

\definehead[exercise][subsection]

\setuphead
  [exercise]
  [number=yes,
   style=,
   numberstyle=\ss\bf\tfb,
   textstyle=\ss,
   alternative=margin,
   numbercommand={\determineheadnumber[subsection]\currentheadnumber\gobbleoneargument}]

\setuphead
  [chapter]
  [number=no,
   style=,
   textstyle=\tfb\ss]

\setuphead
  [section]
  [number=no,
   style=,
   textstyle=\tf\ss]

\definemeasure[headerimageheight][\dimexpr(\headerheight-2ex)]

\definelayer[sabel]

\startsetups[sabel]

\setlayerframed[sabel][x=0mm,y=0mm][frame=off,offset=none]
{%
   \setupTABLE[start][style=\tfx]
   \setupTABLE[r][each][height=0.85\baselineskip]
   \setupTABLE[column][1][align=right, width=.4\textwidth, loffset=-2pt, toffset=-5pt]%
   \setupTABLE[column][2][align={center,lohi}, width=.2\textwidth]%
   \setupTABLE[column][3][align=left, width=.4\textwidth, roffset=-2pt, toffset=-5pt]%
   \bTABLE[frame=off]
   \bTR \bTD Sabel-GBS       \eTD \bTD [nr=5] \dontleavehmode\externalfigure[assets/sabel.png][maxheight=\measure{headerimageheight}, frame=off]\dontleavehmode \eTD \bTD Schuljahr 2025/26 \eTD \eTR
   \bTR \bTD Mathematik      \eTD                                                                                                                                     \bTD Klasse M3A        \eTD \eTR
   \bTR \bTD B.~Donner       \eTD                                                                                                                                     \bTD                   \eTD \eTR
   \bTR \bTD \strut          \eTD                                                                                                                                     \bTD                   \eTD \eTR
   \bTR \bTD 1. Schulaufgabe \eTD                                                                                                                                     \bTD 17.11.2025        \eTD \eTR
   \eTABLE}

\stopsetups

\setupbackgrounds[header][text][bottomframe=on,setups=sabel,background=sabel]
\setuppagenumbering[location={footer,middle},style=\bfc]

\setupfootertexts
  [BE: 67 Arbeitszeit: 60min][Viel Erfolg!]
  [Viel Erfolg!][BE: 67 Arbeitszeit: 60min]


\definedescription[example][
  headstyle=\it, style=, align=left, location=top,
  width=broad, margin=1cm, indenting=never, text=\hskip-1cm]

\definemathcommand [kgV] [nolop] {\mfunction{kg\kern-.15em\hbox{V}}}

\setdigitmode 3

\setupitemgroup[itemize][1][n,repeat][width=1em]
\setupitemgroup[itemize][2][n,repeat][width=2em]
\setupitemgroup[itemize][3][n,repeat][width=3em]
\setupitemgroup[itemize][4][n,repeat][width=4em]

\startMPinclusions
  input "include/metapost/mymacros.mp"
\stopMPinclusions

\starttext

Hinweis: Achten Sie auf eine übersichtliche und schöne äußere Form, korrekte
mathematische Schreibweise und geben Sie bei Textaufgaben kurze Antworten an.
Gestalten Sie saubere und übersichtliche Zeichnungen in aussagekräftiger Größe.

\startitemize
\item Gegeben ist die reelle Funktion \math{g} mit \math{g(x) = (2x +1) e^{-\digits 0.5 x}}
auf ihrer maximalen Definitionsmenge \math{D_g={\blackboard R}}. Der Graph der
Funktion \math{g} wird mit \math{G_g} bezeichnet. Runden Sie Ihre Ergebnisse
gegebenenfalls auf zwei Nachkommastellen.
\startitemize
\item Untersuchen Sie \math{G_g} auf mögliche Schnittpunkte mit den
Koordinatenachsen. Geben Sie gegebenenfalls die Koordinaten der Schnittpunkte an. (BE~3) 
\item Untersuchen Sie das Verhalten der Funktionswerte von \math{f} für
\math{x \to -\infty} und \math{x \to +\infty}. Ermitteln Sie die Gleichung aller
Asymptoten von \math{G_g}. (BE~3)
\item Bestimmen Sie die Art und Lage der Extrempunkte von \math{G_g}. (BE~9)\\
[mögliches Zwischenergebnis \math{g'(x) = (\digits 1.5 -x)e^{-\digits 0.5 x}}]
\item Zeichnen Sie den Graphen \math{G_g} unter Verwendung aller bisherigen
Ergebnisse und weiterer geeigneter Funktionswerte sowie alle Asymptoten von
\math{G_g} mindestens im Bereich \math{-2 \leq x \leq 10} in ein kartesisches
Koordinatensystem. (BE~7)
\item Weisen Sie nach, wieviele Wendepunkte \math{G_g} besitzt und
geben Sie diese gegebenenfalls an. (BE~3)
\item Zeigen Sie, dass die Funktion \math{G(x) = (-4x-10)e^{-\digits 0.5 x}}
eine Stammfunktion von \math{g} ist. (BE~3)
\item Bestimmen Sie den Wert von \math{\int_{-1}^{0} g(x) dx} und interpretieren
Sie das Ergebnis geometrisch. (BE~4)
\stopitemize


\item Gegeben ist die reelle Funktion \math{f} mit \math{f(x)=\ln \frac{x^2+2}{x}}
auf ihrer maximalen Definitionsmenge \math{D_f}. Der Graph der Funktion
\math{f} wird mit \math{G_f} bezeichnet. Runden Sie Ihre Ergebnisse
gegebenenfalls auf zwei Nachkommastellen.
\startitemize
\item Geben Sie die Definitionsmenge \math{D_f} an. (BE~4)
\item Untersuchen Sie die Funktion \math{f} auf mögliche Nullstellen. (BE~4)
\item Bestimmen Sie Art und die Lage der Extempunkte von  \math{G_f}. (BE~11)\\
[mögliches Zwischenergebnis \math{f'(x) = \frac{x^2 - 2}{(x^2 + 2)x}}]\\
[mögliches Zwischenergebnis \math{f''(x) = \frac{-x^4 + 8x^2 + 4}{(x^2 + 2)^2 x^2}}]
\item Untersuchen Sie das Verhalten der Funktionswerte an den Rändern des
Definitionsbereichs. Ermitteln Sie die Gleichung aller Asymptoten von
\math{G_f}. (BE~3)
\item Weisen Sie nach, wieviele Wendepunkte \math{G_f} besitzt und
geben Sie diese gegebenenfalls an. (BE~7)
\item Zeichnen Sie den Graphen \math{G_f} unter Verwendung aller bisherigen
Ergebnisse und weiterer geeigneter Funktionswerte sowie alle Asymptoten von
\math{G_f} mindestens im Bereich \math{-\digits 0.5 \leq x \leq 8} in ein
kartesisches Koordinatensystem. (BE~6)
\stopitemize
 
\stopitemize
\stoptext
